\section*{Erd\H{o}s Problem 559}

\paragraph{Statement and result.}
Does every \(n\)-vertex graph of maximum degree at most a fixed
\(d\) have size Ramsey number \(O_d(n)\)? No, already for maximum
degree three. We reconstruct Tikhomirov's random-component method
with a different, simpler sequence of parameters. It gives infinitely
many orders \(n\) and graphs \(F\) of maximum degree three such that
\[
 \widehat R(F)>\frac{Dn}{100},\qquad
 D=\exp(\Omega(\sqrt{\log n})). \tag{1}
\]
Size Ramsey number counts edges of a host whose every red-blue
colouring contains a monochromatic ordinary copy of \(F\).

\paragraph{A random rooted component.}
Let \(U_k\) consist of the full rooted binary tree of depth \(k\),
with a uniformly random Hamilton cycle added on its \(L=2^k\)
leaves. It has \(s=2L-1\) vertices and maximum degree three.
All its vertices have distance at most \(k\) from its root.

For a fixed graph \(J\) of maximum degree at most \(D\), and a
specified root image \(v\), the probability that \(U_k\) embeds
with root at \(v\) is at most
\[
 Q:=\frac{D^{3L}}{(L-1)!}. \tag{2}
\]
Indeed, there are at most \(D^{2L-2}\) ways to map the tree edges
starting at \(v\), a bound also valid for injective embeddings.
For each fixed embedding, generate the random leaf cycle by fixing
one starting leaf and choosing a uniform ordering of the others.
At each successive choice there are at most \(D\) acceptable next
images. Ignoring the final closing edge gives probability at most
\(D^{L-1}/(L-1)!\). A union bound proves (2).

\paragraph{Choose components with few possible common root images.}
Let \(\ell\) tend to infinity through integers and put
\[
 k=10\ell,\quad D=2^\ell,\quad L=2^k,\quad
 h=2^{k^2},\quad B=D^{k+1},\quad r=DB=D^{k+2}.
\]
Choose \(h\) independent copies \(U_1,\ldots,U_h\) of \(U_k\).
With positive probability they have the following property:
\[
 \begin{gathered}
 \text{In every graph of maximum degree at most }D,\\
 \text{each vertex is a possible root image for fewer than }r
 \text{ of the indexed components.}
 \end{gathered} \tag{3}
\]

To prove this simultaneous assertion, all rooted embeddings in
question lie in a radius-\(k\) ball, whose order is at most
\(1+D+\cdots+D^k\leq B\).
Relabel that ball on \([B]\), with root label 1, adding isolated
vertices if necessary. The number of possible such labelled graphs
of maximum degree at most \(D\) is at most \(B^{DB}=B^r\).
For example, encode each vertex's sorted neighbour list in \(D\)
slots, padding with its own label; this gives an injection into
arrays with \(DB\) entries from \([B]\).

For each choice of \(r\) indexed components and one such graph,
independence and (2) bound the probability that all embed at its
root by \(Q^r\). Thus the probability that (3) fails is at most
\[
 \binom hr B^r Q^r
 \leq\left(\frac{ehQ}{D}\right)^r. \tag{4}
\]
For sufficiently large \(\ell\), the expression in parentheses is
less than one. To check this rather than assume it, use
\(\log((L-1)!)\geq(L-1)(\log(L-1)-1)\). Its logarithm is at most
\[
 1+k^2\log2-\ell\log2+3L\ell\log2
 -(L-1)(\log(L-1)-1).
\]
Dividing by \(\ell L\), this tends to
\((3-10)\log2<0\). Hence (4) is strictly less than one eventually,
which proves existence of a fixed choice satisfying (3).
Let \(F\) be its disjoint union, with \(n=hs\) vertices.

\paragraph{Colour every host with at most \(Dn/100\) edges.}
Fix any such host \(G\), and write \(M=Dn/100\).
Let \(W\) be its vertices of degree at least \(D+1\).
There are fewer than \(2M/D=n/50\) of them.
Delete all edges incident to \(W\), obtaining \(J\) of maximum
degree at most \(D\). There are at most \(2M\) non-isolated vertices
in \(J\). For each \(i\), let \(A_i\) be the set of possible root
images for an embedding of \(U_i\) into \(J\).
Property (3) implies
\[
 \sum_{i=1}^h|A_i|\leq2Mr.
\]
Choose \(i_0\) with \(|A_{i_0}|\leq2Mr/h\).
Colour red all edges incident to \(W\), and all remaining edges
incident to \(A_{i_0}\). Colour all other edges blue.
The blue graph contains no \(U_{i_0}\): any such embedding would
place its root in \(A_{i_0}\), where no blue edge remains.
It consequently contains no \(F\).

There are at most
\[
 D|A_{i_0}|\leq\frac{2MDr}{h}
 =\frac{D^2sr}{50}<\frac h2 \tag{5}
\]
red edges which are not incident to \(W\), for sufficiently large
\(\ell\). Explicitly, \(s<2^{k+1}\), so the binary logarithm of
the expression before the last inequality is at most
\((k+4)\ell+k+1\), whereas \(\log_2h=k^2=100\ell^2\).

Suppose a red copy of \(F\) existed. By (5), fewer than \(h/2\)
of its disjoint components could use an edge not incident to \(W\).
At least \(h/2\) components would therefore have every edge incident
to \(W\). In each of these components the leaf cycle of length
\(L\) needs at least \(L/2\) vertices of \(W\): every vertex
covers at most two of its \(L\) edges.
Disjointness would force
\[
 |W|\geq hL/4>n/8,
\]
contradicting \(|W|<n/50\). There is no red copy either.
Since the host was arbitrary, this proves (1).

\paragraph{Growth and scope.}
Here \(n=2^{100\ell^2}(2^{10\ell+1}-1)\), so
\(\log n=\Theta(\ell^2)\), whereas \(\log D=\ell\log2\).
In particular \(\widehat R(F)/n>D/100\to\infty\), contradicting
any universal linear bound for degree-three target graphs.
An infinite sequence of orders is sufficient for this disproof;
the stronger assertion for every large order is not required here.

\paragraph{Attribution.}
The rooted binary-tree components, random leaf cycles and two-stage
host colouring are reconstructed from K.~Tikhomirov,
\emph{On bounded degree graphs with large size-Ramsey numbers},
Section 2: \url{https://arxiv.org/abs/2210.05818}.
The parameter selection and all probabilistic and counting steps
used above are supplied explicitly; no external size-Ramsey theorem
is assumed. The original disproof is due to R\"odl and Szemer\'edi.
Current statement: \url{https://www.erdosproblems.com/559},
checked 24 September 2026. No novelty is claimed.

\paragraph{Completion estimate.}
The full negative answer follows from a proved superlinear lower bound
for an infinite family of graphs of maximum degree three.
\textbf{COMPLETION: 100\%}
